\section{Motivation}%0.75 pages
\label{sec:motivation}

\subsection{Operator Fusion Opportunities}
\label{sec:fusion}

% 可做两部分实验，以GEMM为例，dynamic shape下，一个是单独算子，另一个是融合算子，单独算子时triton/tilelang能够比得过cublas/cudnn的比例较低；但融合算子对比cutlass等NVIDIA实现的融合算子，很大程度上会显示出优势？

% 从而服务于contribution 1-2，算子融合的必要性，以及triton/tilelang的模板保留了高性能和灵活性


\subsection{Variation in Optimal Configurations}
\label{sec:static}

% 以基于GEMM的融合算子为例，可给出热力图，证明不同形状下最优配置是有差别的，没有one-fits-all，服务于contribution 3

%\subsection{Tile Configuration Trade-offs}
%\label{sec:atuo-fusion}
%
%% 



\subsection{Bottleneck Shifts across Shapes}
\label{sec:bottleneck}

% 以GEMM和卷积算子为例，各画一个roofline model？不同shape下不同configuration的差别，证明比如随着shape变化，更趋向于memory-bound或compute-bound

